\chapter{Appendix 1}

\section{Change of measure for the Black-Scholes model}
\label{sec:change_of_measure}

Consider the discounted stock price
\begin{equation}
    \tilde S_t := \frac {S_t}{B_t} = e^{-rt} S_t
\end{equation}
and apply Itô's lemma to get a new SDE in terms of the discounted stock price
\begin{equation}
    \mathrm d \tilde S_t = \tilde S_t[(\mu-r)\mathrm dt + \sigma \mathrm dW_t^\mathbb P]
\end{equation}
which means that under the physical measure $\mathbb P$, $\tilde S_t$ has nonzero drift unless $\mu = r$, and hence not a martingale. So to enforce no arbitrage in our pricing model, we need a new equivalent probability measure, name it $\mathbb Q$, under which $\tilde S_t$ is a martingale.

Let's define the quantity
\begin{equation}
    c := \frac{\mu -r}{\sigma}
\end{equation}
which measures how much excess drift per unit volatility the stock earns under the physical measure $\mathbb P$. Then we can define the Radon-Nikodym derivative
\begin{equation}
    \frac{\mathrm d\mathbb Q}{\mathrm d \mathbb P}\bigg|_{\mathcal F_t} = Z_t := \exp \left(-cW_t^\mathbb P- \frac 12 c^2 t\right)
\end{equation}
which satisfies Novikov's condition as $c$ is constant so $\mathbb Q$ is well defined. The new process
\begin{equation}
    W_t^\mathbb Q := W_t^\mathbb P + ct
\end{equation}
is a Brownian motion under $\mathbb Q$ by Girsanov's theorem. If we substitute this into the original SDE, we obtain
\begin{align}
    \frac{\mathrm dS_t}{S_t} &= \mu \mathrm dt + \sigma(W_t^\mathbb Q - c\mathrm dt) \\
    &= (\mu - \sigma c) \mathrm dt + \sigma \mathrm dW_t^\mathbb Q \\
    &= r \mathrm dt + \sigma \mathrm dW_t^\mathbb Q
\end{align}

So now under $\mathbb Q$, we have $\mathrm d\tilde S_t = \sigma \tilde S_t \mathrm d W_t ^\mathbb Q$ which has zero drift, and hence $\tilde S_t$ is a $\mathbb Q$-martingale.

Intuitively, this change of measure reweights possible stock paths so that there is no more risk premium, and all assets earn the risk-free rate in expectation after discounting.